\documentclass[11pt]{amsart}
\usepackage{calc,amssymb,amsmath,ulem,stmaryrd,fullpage}
\usepackage{enumerate}
\usepackage{alltt}
\RequirePackage[dvipsnames,usenames]{color}
\let\mffont=\sf
\normalem
%\input{kmacros3.sty}
\input{xy}
\xyoption{all}
\usepackage{hyperref}
\usepackage{graphicx}
\usepackage[all,cmtip]{xy}
\usepackage{verbatim}
\usepackage[left=0.95in,top=0.95in,right=0.95in,bottom=0.95in]{geometry}
\newcommand{\tauCohomology}{T}
\newcommand{\FCohomology}{S}


\theoremstyle{theorem}
%\newtheorem*{mainthm*}{Main Theorem}

\newcommand{\note}[1]{\marginpar{\sffamily\tiny #1}}

\def\todo#1{\textcolor{Mahogany}%
{\sffamily\footnotesize\newline{\color{Mahogany}\fbox{\parbox{\textwidth-15pt}{#1}}}\newline}}

\renewcommand{\O}{\mathcal O}
\newcommand{\ie}{{\itshape i.e.} }
\newcommand{\roundup}[1]{\lceil #1 \rceil}
\newcommand{\rounddown}[1]{\lfloor #1 \rfloor}
\renewcommand{\to}[1][]{\xrightarrow{\ #1\ }}
\newcommand{\onto}[1][]{\protect{\xrightarrow{\ #1\ }\hspace{-0.8em}\rightarrow}}
\newcommand{\into}[1][]{\lhook \joinrel \xrightarrow{\ #1\ }}
%\newcommand{DBDual}[1]{{\underline{\omega}_{#1}}}
\newcommand{\DBDual}[1]{{{\uuline \omega}{}^{\mydot}_{#1}}}
\newcommand{\Tt}{{\mathfrak{T}}}
%\newcommand{\bn}{{\mathfrak{n}} }
\newcommand{\sol}[1]{ \vskip 6pt{\color{blue} {\bf Solution:} #1} \vskip 6pt}

\newcommand{\bR}{{\mathbb{R}}}
\newcommand{\va}{{\vec{a}}}
\newcommand{\vb}{{\vec{b}}}
\newcommand{\vzero}{{\vec{0}}}
\newcommand{\vi}{{\vec{i}}}
\newcommand{\vj}{{\vec{j}}}
\newcommand{\vk}{{\vec{k}}}
\newcommand{\vh}{{\vec{h}}}
\newcommand{\vr}{{\vec{r}}}
\newcommand{\vu}{{\vec{u}}}
\newcommand{\vv}{{\vec{v}}}
\newcommand{\vw}{{\vec{w}}}
\newcommand{\vx}{{\vec{x}}}
\newcommand{\vy}{{\vec{y}}}
\newcommand{\vz}{{\vec{z}}}
\newcommand{\vT}{{\vec{T}}}
\newcommand{\vN}{{\vec{N}}}
\newcommand{\vF}{{\vec{F}}}
\newcommand{\vG}{{\vec{G}}}
\newcommand{\bF}{\mathbb{F}}
\newcommand{\bQ}{\mathbb{Q}}
\newcommand{\bZ}{\mathbb{Z}}
\newcommand{\Inn}{\mathrm{Inn}}
\newcommand{\Aut}{\mathrm{Aut}}


\begin{document}
\title{Worksheet \#1 -- Math 6310\\Fall 2019}
\author{Due September 4th, 2019}
\maketitle
Let's play around with universal properties.

In general category theory, an \emph{epimorphism} between objects is something that is right cancellative.  In other words, $f$ is an epimorphism, if we have
\[
\xymatrix{
A \ar[r]^f & B \ar@/_1pc/[r]_h \ar@/^1pc/[r]^g & C
}
\]
a diagram, and if we have $g \circ f = h \circ f$, then $g = h$.  (These could be groups and group homomorphisms, or sets and ordinary functions, or ...).
\vskip 9pt
\noindent
{\bf 1.}  Suppose that $A, B, C$ above are sets and $f,g,h$ as above are just functions between sets.  Show that if $f$ is surjective then $f$ is an epimorphism.
\vskip 170pt
\noindent
{\bf 2.}  Prove the converse to {\bf 1.}, if $f$ is an epimorphism, then $f$ is surjective.
\newpage
Now we switch to the category of groups.  In other words suppose we have $f : A \to B$ a homomorphism of groups.  We declare $f$ to be an \emph{epi} if $g \circ f = h \circ f$ implies $g = h$.
\[
\xymatrix{
A \ar[r]^f & B \ar@/_1pc/[r]_h \ar@/^1pc/[r]^g & C
}
\]
\vskip 9pt
\noindent
{\bf 3.}  In the category of \emph{Abelian groups} (ie, assume, $A,B,C$ are Abelian), show that $f$ is an epi if and only if $f$ is surjective.  (This is still true without the Abelian assumption, but it's a bit harder, and may not fit in the space provided.)
\vskip 3pt
\emph{Hint:}  If $f$ is surjective, it should easily be an epi (from the previous page).
\vskip 175pt
Now let's think about the category of monoids.  Consider the inclusion map $f : (\bZ_{\geq 0}, +, 0) \to (\bZ, + , 0)$.  Obviously this map is not surjective, however see the next exercise.
\vskip 9pt
\noindent
{\bf 4.}  Show that $f$ is an epi in the category of monoids.
\vskip 3pt
\emph{Hint:} Is it true for monoids that $g(-n) = -g(n)$?
\end{document}

